\begin{pastexam}[label={px:ntu-115-dsa-24}]{臺大資工所 115 學年度 第 24 題}
（單選）

(4 points) Let $G=(V,E)$ be a directed graph represented using adjacency lists without edge weights. A standard Breadth-First Search (BFS) is run from a source vertex $s$, producing distance labels $d[v]$ for all reachable vertices $v$. Which of the following statements is always true?
\begin{enumerate}[label=(\Alph*)]
\item For every edge $(u,v)\in E$, we must have $d[v]\le d[u]+1$.
\item If $d[v]=d[u]+1$, then $(u,v)$ must be a tree edge in the BFS tree.
\item In a directed graph, BFS always produces a shortest-path tree that is also a topological ordering of the reachable subgraph.
\item If the graph is represented by an adjacency matrix instead of adjacency lists, BFS may compute incorrect distances due to repeated edge relaxations.
\item If BFS discovers vertices in the order $v_1,v_2,\ldots,v_k$, then this order is a shortest-path ordering for all pairs of vertices.
\end{enumerate}
\end{pastexam}

\begin{pastsoln}
答案是 (A)。
\begin{itemize}
\item (A) 對。 資結書第~21.3.4~節的事實 (a)：每條邊 $(u, v)$ 都有 $\delta(s, v) \le
\delta(s, u) + 1$；資結書定理~21.4 說 BFS 結束時 $d = \delta$，所以 $d[v] \le d[u] + 1$。
題目只替可到達的頂點定義 $d$；$u$ 到不了時照資結書的約定 $d[u] = \infty$，不等式自動成立（$u$ 可到
達時 $v$ 也可到達，不會出現 $d[v]$ 沒有定義的情形）。
\item (B) 錯。 反例：邊 $s \to a$、$s \to b$、$a \to c$、$b \to c$，串列依字母排。 BFS 依序發現
$a$、$b$（$d = 1$），拿出 $a$ 時發現 $c$，$\pi[c] = a$、$d[c] = 2$。 $(b, c)$ 也滿足 $d[c] =
d[b] + 1$，卻不是樹邊：每個頂點只有一個 $\pi$，由誰先發現決定。
\item (C) 錯。 可到達的部分可能有環，例如 $s \to a$、$a \to s$，根本沒有拓樸順序（第~22~章）。
就算沒有環，BFS 的順序也不一定是拓樸順序：邊 $s \to a$、$s \to b$、$b \to a$，發現順序是
$s, a, b$，但邊 $b \to a$ 從後面指向前面。
\item (D) 錯。 表示法只影響時間（矩陣是 $\Theta(n^2)$）。 BFS 的 $d[v]$ 在 $v$ 第一次被發現時設定一
次，之後不再修改（資結書虛擬碼~21.1 第 11--15 行只處理白色頂點），沒有「重複放鬆」這回事，
算出來的距離和用串列時一樣。 放鬆是第~25~章加權最短路徑的操作。
\item (E) 錯。 發現順序只保證「從 $s$ 出發的距離」不減（資結書第~21.3.4~節的 (c)），
和其他頂點之間的距離無關。 用 (C) 的第二個例子：發現順序是 $s, a, b$，但從 $a$ 根本到不了 $b$，
從 $b$ 到 $a$ 卻只要一步。
\end{itemize}
\end{pastsoln}
